\begin{description}[leftmargin=0pt,itemsep=2pt]
\item[Does lower prediction loss correspond to better features?] No. Within a run the held-out loss is minimal at step 700,
when accuracy and rank are still at random-encoder level, and then \emph{rises} while accuracy improves by 13 points.
% ($\rho{=}{+}0.65$ between loss and validation accuracy after step 500). 
The 500-step model has the lowest loss of the Part~A
checkpoints (0.078) and random-level features. 
After the initial drop in the training loss, the loss starts increasing due to the target encoder shifting away from 
position-dependent features towards semnatic concepts making it harder for the context encoder and the predictor to learn.
This forces the context encoder to learn useful contextual and semantic features from the availabel context features, 
for the predictor to predict the hidden target features. This is supported by the increasing gaps between the validation 
heldout latent loss, context-free loss and other-image context loss.
% Across variants the lowest loss belongs to
% scattered (0.113), the least noise-robust run (17.0 at $\sigma{=}0.15$). The early drop is learnt from positions: the loss with another image's
% context falls exactly like the real-context loss, and a content-free per-position predictor beats both until step 250
% (Sec.~\ref{sec:h2}). The effective rank, by contrast, tracks clean accuracy.
%  ($\rho{=}{+}0.92$).

\item[Does the change improve clean classification, occlusion robustness, or both?] The geometry change 
(1\_large) is the most
noise-robust geometry ($27.5{\pm}2.3$ vs $22.2{\pm}1.0$ at $\sigma{=}0.15$, 5 seeds each) and improves occlusion robustness
(+2.3 / +2.8 points at 25 / 40\%) with a small clean gain (+0.9); it slows the noise failure but does not repair it (it stays
at random-encoder level). Pretraining with context noise repairs the noise failure (+19.7 at $\sigma{=}0.15$) at a small clean cost ($-$1.4 with the
selected probe) and leaves occlusion unchanged. 
% Neither closes the gap to the oracle probe at 40\% occlusion (38.4\%), which
% is the part of the drop due to a feature shift rather than to missing evidence.

\item[Are gains caused by mask geometry, or by a change in visible/target counts?] By geometry. All H1 candidates use exactly
20 context and 32 target tokens per image (scattered: 32--36), and the scattered
ablation with the same counts is the least noise-robust run (17.0 vs $27.5{\pm}2.3$ for 1\_large at $\sigma{=}0.15$), with
halved image-specific feature variance and a held-out loss close to a position-only prediction. On occlusion the separation
is weaker (scattered, one seed, matches 1\_large at 25\% and is 1.1 points lower at 40\%), so the occlusion gain is only
partly due to contiguity. Within ``geometry'',
contiguity and block scale are confounded (1\_large has one large contiguous block, scattered many small ones).

\item[Could extra parameters, different training exposure, or a probe change explain the result?] No for parameters: every
run has 1{,}412{,}864 trainable parameters. No for exposure: same 15k steps, same seeded data order and images, and the
predictor sees fewer target tokens than the baseline (32 vs 45). No for the probe: identical protocol, $(\mathrm{lr},\mathrm{wd})$
selected on validation; the same ranking holds under the fixed probe setting (e.g.\ 1\_large 51.0 clean / 31.5 at 40\%
occlusion, baseline 48.8 / 28.2). According to the paper, increasing the depth of the predictor can help with improved 
downstream performance, but significant changes can only be seen when training on larger image sizes and larger dataset size.
% For context noise the extra exposure \emph{is} the mechanism; its noise is drawn from a
% range chosen knowing the strongest test level ($\sigma_{\max}{=}0.2>0.15$), with seeds disjoint from the test noise.

\item[Could apparent feature stability come from collapsed or uninformative representations?] Yes, and this is observed in 
the training loss graph where the loss initially drops and the heldout losses (context free, and other image's context) drop
converge. At this point, the features are collapsed since they context encoder and predictor depend on the target positions
to predict the target features and the target encoder's weights are still not fully changed from its random weights. But as
the training proceeds, the target encoder shifts away and produces high image and patch variant target features, forcing the 
context encoder and predictor to move beyond just depending on the target position embeddings and focus on contextual and 
semantic concepts in the image. 

% Yes, and the random encoder is
% the clearest case: its features are the most stable under noise (clean--noisy cosine 0.997) and the least informative
% (35\% clean). The baseline's stability falls as its rank grows (cosine 0.71 at 15k). Stability is only meaningful together
% with rank and accuracy: context-noise pretraining is stable (0.95) \emph{and} informative (rank 57.7, 47.5\% clean). Scattered
% masking shows the opposite pitfall: high pairwise cosine (0.60) from a halved image-specific variance without loss of rank.

\item[Which single follow-up experiment would you prioritise with one more week?] To evaluate the pretrained encoders on low-level
vision tasks like counting and segmentation and study the performance. To also scale the image size and dataset size to large datasets,
and scale the encoder and predictor sizes, to study how it improves downstream performance and in which scenarios its features can collapse.
% Evaluate context-noise pretraining (3 seeds,
% combined with and without 1\_large) on corruptions it never saw, at a maximum training noise \emph{below} the test level
% ($\sigma_{\max}{=}0.1$) and on other families (blur, JPEG, impulse noise, contrast; the standard CIFAR-10-C set). If robustness
% transfers, the diagnosis ``missing invariance pressure'' is confirmed and the fix is a method; if it only helps at the trained
% $\sigma$, it is test-condition matching. This single experiment decides which of the two the noise result is.
\end{description}

\subsection*{Final research question: what limits robust representation learning here?}
% Pretraining with only the initial training budget of 500 steps limits the encoder from learning any useful features and results in 
% poor downstream performance. 
\begin{itemize}
\item \textbf{Image size} The small image size $32\times32$ limits different masking patterns and strategies when experimenting with larger target blocks.
\item \textbf{Not feature collapse.} Every run ends with effective rank 57--66, top-1 singular-value share $\le$0.16 and growing
feature spread; the one contraction we found (scattered) leaves the rank intact.
\item \textbf{Not context size or target overlap.} Lowering the context to 16.8 tokens and capping the overlap between the
4 small blocks change accuracy by at most 1.5 points.
\item \textbf{Mask geometry decides what the encoder relies on.} Large targets far from the context keep image-level
features and the best noise robustness of any geometry (1\_large, +5.3 points at $\sigma{=}0.15$); small targets next to the
context produce local, position-driven features and the worst (scattered, $-$5.2). Geometry alone does not make the encoder
noise-invariant, and adds $\approx$3 points under 40\% occlusion.
\item \textbf{Limited training is the bottleneck only at the assignment budget.} At 500 steps the encoder equals a random encoder;
longer training improves clean accuracy but \emph{worsens} noise robustness.
\item \textbf{Best-supported explanation: the objective has no invariance pressure, and a clean-only probe protocol exposes it.}
I-JEPA's target is the clean image's own features, so the encoder grows features for fine local structure that noise
perturbs; the class information survives (39.6\% with a probe trained on mixed clean and noisy images) but moves in feature space. Putting the
perturbation on the context side alone recovers it (41.9\% with the clean probe). For occlusion the largest part of the drop
is missing evidence (the supervised reference loses the same fraction).
\item \textbf{Next experiment:} Scaling the image and dataset sizes, scaling the encoder and predictor networks, 
and measure the feature robustness against noise, occusion. To also evaluate the features on low-level vision tasks.
\end{itemize}
